\clearpage
\section{유한체 위의 원분다항식}
\label{sec:finite-field-cyclotomy}

\begin{learninggoal}
유한체의 상대 프로베니우스가 단위근에 지수 $q$로 작용함을 이용한다. 법 $n$에서 $q$의 곱셈위수로 원분다항식의 기약인수 차수와 분해체를 계산한다.
\end{learninggoal}

$q$를 소수거듭제곱이라 하고 $\gcd(q,n)=1$이라 하자. 대수적 폐포 $\overline{\F}_q$에서 원시 $n$차 단위근 $\zeta_n$을 택한다. 유한체의 프로베니우스는
\[
  \Frob_q:x\longmapsto x^q
\]
이고 단위근 위에서는
\[
  \Frob_q(\zeta_n)=\zeta_n^q
\]
로 작용한다.

\begin{theorem}[유한체 위 원분확장의 차수]
\label{thm:finite-field-cyclotomic-degree}
$d=\operatorname{ord}_n(q)$를 $q$의 $(\Z/n\Z)^\times$에서의 곱셈위수라 하자. 그러면
\[
  \boxed{[\F_q(\zeta_n):\F_q]=d}
\]
이고
\[
  \Gal(\F_q(\zeta_n)/\F_q)
  =\langle\Frob_q\rangle
  \cong\langle q\rangle
  \le(\Z/n\Z)^\times.
\]
특히
\[
  \F_q(\zeta_n)=\F_{q^d}.
\]
\end{theorem}

\begin{proof}
$\Frob_q^e$가 $\zeta_n$을 고정할 필요충분조건은
\[
  \zeta_n^{q^e}=\zeta_n
  \quad\Longleftrightarrow\quad
  q^e\equiv1\pmod n
\]
이다. 따라서 $\zeta_n$의 프로베니우스 궤도길이는 정확히 $d$이다. 유한체 위 최소다항식의 차수는 프로베니우스 궤도길이와 같으므로 확장차수도 $d$이다. 유한체 확장의 갈루아군은 상대 프로베니우스로 생성되는 차수 $d$의 순환군이다.
\end{proof}

\begin{corollary}[원분다항식의 인수 차수]
\label{cor:finite-field-cyclotomic-factorization}
$d=\operatorname{ord}_n(q)$라 하자. 그러면 $\Phi_n$은 $\F_q[X]$에서
\[
  \boxed{\frac{\varphi(n)}d}
\]
개의 서로 다른 일계수 기약다항식의 곱으로 분해되고, 각 기약인수의 차수는 모두 $d$이다.
\end{corollary}

\begin{proof}
$\Phi_n$의 근은 원시 $n$차 단위근들이다. 각 원시근 $\zeta_n^a$의 프로베니우스 궤도길이도
\[
  \min\{e>0:q^e a\equiv a\pmod n\}=d
\]
이다. 기약인수 하나의 근 집합은 하나의 프로베니우스 궤도이므로 모든 인수의 차수가 $d$이고, 전체 근의 수 $\varphi(n)$을 $d$로 나누면 인수의 개수를 얻는다.
\end{proof}

\begin{corollary}[기약성 판정]
\label{cor:finite-field-cyclotomic-irreducible}
$\Phi_n$이 $\F_q[X]$에서 기약일 필요충분조건은
\[
  \operatorname{ord}_n(q)=\varphi(n)
\]
이다. 이는 $q$의 합동류가 $(\Z/n\Z)^\times$ 전체를 생성한다는 뜻이다.
\end{corollary}

\begin{example}[$\Phi_5$ over $\F_2$]
법 $5$에서
\[
  2,\ 2^2=4,\ 2^3=3,\ 2^4=1
\]
이므로 $\operatorname{ord}_5(2)=4=\varphi(5)$이다. 따라서
\[
  \Phi_5=X^4+X^3+X^2+X+1
\]
은 $\F_2$에서 기약이고, 원시 $5$차 단위근은 $\F_{16}$에 처음 나타난다.
\end{example}

\begin{example}[$\Phi_7$ over $\F_2$]
$2^3\equiv1\pmod7$이고 더 작은 양의 지수에서는 $1$이 아니므로
\[
  \operatorname{ord}_7(2)=3.
\]
따라서 $\Phi_7$은 두 개의 기약삼차식으로 분해된다. 실제로
\[
  \Phi_7(X)
  =(X^3+X+1)(X^3+X^2+1)
  \quad\text{in }\F_2[X].
\]
원시 $7$차 단위근은 $\F_8^\times$의 원소이며, $|\F_8^\times|=7$이므로 영이 아닌 모든 원소가 $7$차 단위근이다.
\end{example}

\begin{example}[$\Phi_8$ over $\F_3$]
$3^2\equiv1\pmod8$이고 $3\not\equiv1\pmod8$이므로 $d=2$이다. 따라서 $\Phi_8=X^4+1$은 두 기약이차식의 곱이다.
\[
  X^4+1
  =(X^2+X+2)(X^2+2X+2)
  \quad\text{in }\F_3[X].
\]
분해체는 $\F_9$이다.
\end{example}

\begin{example}[$X^5-1$ over $\F_7$]
\[
  X^5-1=(X-1)\Phi_5(X).
\]
법 $5$에서 $7\equiv2$이고 $\operatorname{ord}_5(2)=4$이므로 $\Phi_5$는 $\F_7$에서 기약이다. 따라서 $X^5-1$의 분해체는
\[
  \F_{7^4}
\]
이고 갈루아군은 위수 $4$의 순환군이다.
\end{example}

\subsection{$\Q$와 유한체의 비교}

\begin{center}
\renewcommand{\arraystretch}{1.25}
\begin{tabular}{lll}
\toprule
기준체 & 갈루아군의 이미지 & 원분다항식의 차수 \\
\midrule
$\Q$ & $(\Z/n\Z)^\times$ 전체 & $\varphi(n)$, 항상 기약 \\
$\F_q$ & $\langle q\rangle$ & $\operatorname{ord}_n(q)$인 인수들 \\
\bottomrule
\end{tabular}
\end{center}

$\Q$에서는 모든 단위원 지수 $a$가 자동동형 $\zeta_n\mapsto\zeta_n^a$로 나타난다. 유한체에서는 하나의 프로베니우스 $q$가 생성하는 부분군만 나타난다. 따라서 같은 $\Phi_n$도 $\Q$에서는 기약이지만 $\F_q$에서는 여러 같은 차수의 인수로 갈라질 수 있다.

\begin{remark}
조건 $\gcd(q,n)=1$은 단순한 편의가 아니다. 표수가 $n$을 나누면 $X^n-1$ 자체가 중근을 가지며, 원시 $n$차 단위근이 존재하지 않을 수 있다. 이 경우에는 $n$에서 표수의 거듭제곱 부분을 제거한 뒤 서로 다른 근의 구조를 먼저 살펴야 한다.
\end{remark}

\begin{keyidea}
유한체의 원분 계산은 하나의 합동식으로 압축된다.
\[
  \boxed{
  \zeta_n\in\F_{q^d}
  \quad\Longleftrightarrow\quad
  n\mid(q^d-1).}
\]
가장 작은 그런 $d$가 $\operatorname{ord}_n(q)$이고, 이것이 확장차수·프로베니우스 궤도길이·기약인수 차수를 동시에 준다.
\end{keyidea}

\begin{checkpoint}
\begin{enumerate}[label=(\alph*)]
  \item $\Phi_{13}$이 $\F_2$에서 몇 차 기약인수 몇 개로 분해되는지 구하라.
  \item 원시 $9$차 단위근이 $\F_5$의 몇 차 확장에 처음 나타나는지 구하라.
  \item $\Phi_{15}$이 $\F_2$에서 기약인지 판정하고 인수들의 차수를 구하라.
\end{enumerate}
\end{checkpoint}
